\section{Introduction}
A coordinate sweep sends each individual card to a uniform position, while correlations among cards remain. We study a stronger aggregate question: how small is the sum of fixed powers of all the nonconstant singular values of the full permutation operator? The answer gives a quantitative bound for repeated forward sweeps, including their multiplicities in the regular representation.

Thorp introduced the shuffle in \cite{Thorp1973}. For $n=2^d$ positions, Morris obtained an $O(d^{44})$ mixing bound using evolving sets and a chameleon process \cite{Morris2008}. Montenegro and Tetali sharpened this to $O(d^{29})$ through spectral-profile and evolving-set estimates \cite[Section~6.4]{MontenegroTetali2006}. Morris then developed an entropy method based on card interactions, obtaining $O(\log^4 n)$ for every even deck size \cite{Morris2009} and $O(d^3)$ for dyadic decks \cite{Morris2013}. These bounds count physical shuffles.

The proof here combines partial-deck smoothing with representation multiplicities and an estimate for the overlap of row and column symmetries. Splitting the binary coordinates into two consecutive groups makes one part of a sweep act independently within rows of a rectangular board and the other within columns. The two symmetry decompositions do not commute. Our geometric step estimates their overlap by testing tensor vectors against random subspaces of small dimension, and the resulting induction controls a fixed singular-value moment in the full regular representation.

\subsection{The regular-trace statement}
Index positions by $\mathbb F_2^d$. One physical shuffle sends
\[
 (x_1,x_2,\ldots,x_d)\longmapsto
 (x_2,\ldots,x_d,x_1+\xi_{x_2,\ldots,x_d}),
\]
with independent fair bits $\xi$. Removing the deterministic rotations gives a sweep through the coordinate matchings. The rotation returns to the identity after $d$ steps. Write $Q_n$ for the law of this sweep and for its average in a representation. A permutation maps inputs to outputs; a product applies its rightmost factor first. Different sweeps use independent coins.

Every matching average is an orthogonal projection, so $Q_n$ is a contraction. For $|Y|=(Y^*Y)^{1/2}$, all traces and Schatten norms are unnormalized. The regular trace on the symmetric group is
\[
 \operatorname{Tr}_{\rm reg}|Q_n|^P
 =\sum_{\lambda\vdash n}D_\lambda
       \operatorname{Tr}_{V_\lambda}|Q_n(\lambda)|^P,
\]
where $V_\lambda$ is the irreducible representation of dimension $D_\lambda$. The constant representation contributes exactly one.

\begin{theorem}[Regular-trace smoothing]\label{ten:lowplanes:main}
There is an absolute $P\ge2$ such that for every dyadic $n\ge2$,
\begin{equation}\label{ten:lowplanes:eq:1}
 \operatorname{Tr}_{\rm reg}|Q_n|^P\le1+n^{-10}.
\end{equation}
If $v$ is an integer with $2v\ge P$, then $v$ independent sweeps have total-variation distance at most $\frac12n^{-5}$ from uniform, from every deterministic initial deck. Here $\|\mu-\nu\|_{\rm TV}=\tfrac12\sum_g|\mu(g)-\nu(g)|$.
\end{theorem}
Since one sweep consists of $d$ physical shuffles, the theorem gives an $O(d)$ mixing bound for $n=2^d$ cards. This order is optimal: the support count in Proposition~\ref{ten:support} shows that total-variation distance $1/4$ requires at least $2d-O(1)$ physical shuffles.

The regular trace contains all irreducible multiplicities, not just one copy of each Fourier block. The conclusion concerns powers of the forward sweep. Its proof uses Schatten H\"older and does not identify $(Q_n^*Q_n)^v$ with $(Q_n^v)^*Q_n^v$.

\subsection{Three ranges of representations}
The proof separates three ranges, as shown in Figure~\ref{fig:three-ranges}, because the geometric test has a useful, explicit domain. If $X=\log D_\lambda$, Proposition~\ref{prop:grid-overlap} proves the overlap estimate in a window
\[
 ne^{-A_0\sqrt{\log n}}\le X\le A_1n
\]
for fixed positive $A_0,A_1$. In a balanced $a$ by $b$ board, take row and column Fourier projections selecting real unit carrier vectors, including every multiplicity copy. If their carrier dimensions are $D_R,D_F$, we prove
\[
 D_\lambda\operatorname{Tr}(E_RE_F)
 \le D_RD_F\exp\!\left(\frac{X}{(\log n)^{1/3}}\right).
\]
The test chooses independent, unitarily invariant \emph{complex} $r$-planes in the even and odd tensor alphabets, where $r=\lfloor n^{1/100}\rfloor$. The real assumption concerns the selected carrier vectors, and later permits real spectral decompositions of the sweep factors.

Two analytic estimates prepare the other ranges. Corollary~6.2 of the companion \emph{Routing densities and representation contraction for Thorp sweeps} \cite{OpenAI2026Routing} shows that a forward sweep followed by a reverse sweep, in either orientation, has a fixed power whose density on every ordered $l$-card injection is at most $e^{Cl}$. This handles very large dimensions. For diagrams with a very long first row, we break the sweep into coarse coordinate lines and deform each local operator by its polar decomposition. The two orientations of smoothing control the two ends of the singular-value decomposition. Exact uniformity for one card is preserved under this deformation. Contact counts and representation multiplicities then make the sparse regular-trace contribution at most $n^{-12}$ (Proposition~\ref{prop:sparse-rows}).

The remaining representations lie in the window of the geometric estimate. Orthogonal row and column tags turn the projection overlap into the norm of an invariant tensor. The random-plane tests retain a controlled fraction of the relevant representation. On passing both tests, nearly all slots lie in a small even subspace. Compression to the distinct cell tags then produces an exponential saving. The argument bounds arbitrary tagged vectors; it does not require them to factor across slots.

The Araki--Lieb--Thirring inequality \cite{araki1990,audenaert2007} combines the row and column estimates. An exponent increase by $1+2(\log n)^{-1/4}$ absorbs the geometric error. These increases have bounded product along the halving of bit lengths. The finite-size strict gap provides the base cases.

We use classical branching, Pieri, hook-length and ordinary Schur--Weyl theory \cite{Sagan2001,Stanley1999,VershikOkounkov2005}. The hook support of the even/odd alphabet is the classical Berele--Regev construction \cite{BereleRegev1987}, restated in \cite[Theorem~2.1]{Regev2013}. We derive the needed even/odd realization by applying ordinary Schur--Weyl separately to the two species and inducing, with a sign twist on the odd factors; the proof also tracks the exceptional letters needed outside the hook. The random-plane averaging is an application of unitary invariance and highest-weight dimension estimates. The compact coherent-orbit resolution used in the final compression is the one described by Perelomov \cite[Section~2]{Perelomov1972}; the tag construction and its quantitative trace estimate are proved here.

\begin{figure}[ht]
\centering
\begin{tikzpicture}[node distance=4mm,box/.style={draw,rounded corners,align=center,text width=38mm,minimum height=13mm,font=\small},>=Stealth]
\node[box] (large) {large $\log D_\lambda$\\partial-deck smoothing};
\node[box,right=6mm of large] (sparse) {long first row\\local polar deformation};
\node[box,right=6mm of sparse] (window) {remaining window\\random-plane overlap};
\node[box,below=7mm of sparse,text width=55mm] (trace) {fixed regular trace\\$1+n^{-10}$};
\draw[->] (large.south)--(trace.west);
\draw[->] (sparse.south)--(trace.north);
\draw[->] (window.south)--(trace.east);
\end{tikzpicture}
\caption{The three estimates in the proof of Theorem~\ref{ten:lowplanes:main}. After the annihilated types are removed, the cases are tested successively. The last class lies in the dimension window required by the random-plane estimate.}
\label{fig:three-ranges}
\end{figure}
